\section{实验设置}

\subsection{自由形态器件参数化和物理堆栈}

本研究专注于在$64\times64$平面网格上表示的自由形态二值超表面。这种表示是故意广泛的：它避免将搜索限制在小型解析族中，并以自由形态算子设计实际困难的方式使逆向问题变得困难。这些结构在 fabrication-oriented 图案化的意义上是二值的，数据生成例程施加与预期设计空间一致的对称先验。特别是，结构生成和增强步骤强制执行$C_4+\sigma_x+\sigma_y$对称，也就是一类$D_4$型对称先验。这样定义设计空间，并不是任意裁剪几何自由度，而是为了减少明显的方位冗余、强化面内各向同性，并让$0^{\circ}$--$45^{\circ}$成为具有代表性的基础扇区，使整个超表面或整个方位响应可以通过对称折叠被重建。结果设计空间足够广泛以容纳几个算子族，但又足够稀疏，使得朴素搜索迅速变得浪费。

本工作中所有电磁标注和验证都用对应于嵌入在空气输入和二氧化硅输出介质之间的硅图案层的固定物理堆栈进行。整个代码库中使用的名义几何和材料参数是$500\,\mathrm{nm}$周期、$500\,\mathrm{nm}$图案层厚度、硅作为结构材料、空气作为入射介质、$\mathrm{SiO}_2$作为输出介质。这些选择固定不是因为它们是普遍最优的，而是因为本文关注的是逆向设计表述而非详尽的材料堆栈探索。

\subsection{角度--频率采样和响应表示}

响应场在从$800$到$1300\,\mathrm{nm}$以$50\,\mathrm{nm}$步长的波长网格和从$-40^{\circ}$到$40^{\circ}$以$5^{\circ}$步长的角度网格上采样。结果$11\times17$网格足够密集以编码本研究中针对的广泛传递函数几何，同时对于RCWA标注和重排序保持计算可处理性。保留两个偏振通道，对应于训练数据和任务定义中使用的$T_{pp}$和$T_{ss}$幅度响应。因此，每个样本由一对角度--波长图而非标量分数或一维光谱表示。

这种离散化对可以和不可以声称的内容有实际后果。它支持其区分物理在选定角度和光谱窗口尺度上可见的算子族，包括二阶和四阶微分、低通和高通角度滤波、偏振无关操作和偏振复用。同时，它不声称解析任意尖锐的共振特征或所有可能的相位敏感效应。本文的目标是显示统一的响应空间表述在现实的计算预算内是可行的，而非声称对每个光子可观测量的完全闭合。

\subsection{训练和推理协议}

前向替代模型在该仿真基础数据集上训练，以从候选结构预测完整响应张量。在当前实现中，模型在归一化响应空间中用$L_1$重建目标优化，而物理空间误差在验证期间跟踪。数据增强仅限于在所采用约定下保持光学响应不变的对称保持翻转。然后，条件扩散模型在同一响应张量上训练，连同鼓励采样期间物理一致性和二值化的辅助项。这些选择作为架构细节的重要性不如作为意图指示：生成的结构不仅预期类似于训练样本，而且预期保持为经验证的电磁评估的有意义候选。

推理遵循基于候选的协议而非单次预测协议。对于给定的目标响应，扩散模型产生多个候选结构。前向替代模型用于筛选或近似排序这些候选，之后用RCWA重新评估所选结构。最终排序使用从验证场派生的任务感知响应指标。这种候选级程序使逆向问题的"一对多"特性明确。有用的逆向设计器不应仅返回一个合理结构；它应足够密集地填充可行集，使得生成后物理选择变得有价值。

\subsection{任务、指标和评估范围}

大NA、高透过率、宽带二阶微分器被作为代表性演示，因为它要求高结构化的角度响应，同时在响应空间和图像空间中保持物理可解释性。选择额外任务不是为了扩大应用组合，而是为了测试相同的响应空间表述在目标的受控变形下是否保持连贯。四阶微分改变目标曲率。低通和高通任务改变目标包络。偏振无关和偏振复用任务探测偏振是否自然进入相同表述。时空窗口目标作为更广泛的响应域目标的扩展示例保留。

评估指标反映这种组织。对于逐行算子目标，排序逻辑分离中心抑制、目标形状一致性、边缘行为和有效带宽。对于偏振敏感任务，分析添加联合分数、无源通道抑制和围绕$\pm40^\circ$的通道匹配。在基线评估工具中，最佳情况分数、平均分数、前3分数、超过阈值的成功率、光谱MAE、二值率和候选多样性都显式计算。这些量不是作为任意工程诊断集合报告的。它们对应于本文的核心问题：响应形状是否正确、物理验证是否与名义预测一致、是否保持候选多样性、生成的结构是否保持作为二值超表面的实际可用性。
